\documentclass{article}\usepackage{amsmath,amssymb}
\title{Full-library size and ambient-family pin}\date{16 November 2026}
\begin{document}\maketitle
The command-indexed, free-global and flat-descent feedback statements use a fixed adequate universe. World/input/boundary sorts, root-frame sets, primitive-schema sets and each system's states and menus are small in that universe (or are fibers in a declared larger universe). For a fixed sort, finite formal terms over these sets form a set: they are contained in the countable union over finite tree shapes of finite products of the label and frame sets. No reachability or probability-support restriction is used to obtain this smallness.

The full family category must admit all these small shape-indexed families, with the finite limits used in each actual coefficient diagram. A full library may therefore be uncountable even when every primitive raw row and every selected scope is finite. The proof does not place that library in a category admitting only finite or only countable shape collections without an additional countability hypothesis. If all relevant sorts, frames and names are countable, the finite-term shape sets are countable, but that is an optional independently checked specialization.

Each raw selected scope uses finitely many primitive actors and its actual context/full-content coefficient construction takes a finite limit. Each formal terminal tree has finitely many coefficient factors. Forming the set of all formal trees or schema names does not take an infinite tensor product or an infinite coefficient limit. Consequently the size pin supplies no accumulated infinite-content semantics and no closure theorem for a finite/countable-only source presentation. This is the explicit ambient assumption recorded by the independent review SHA256 \texttt{f40b67240ecf42592b9a5fa7b153096a5d18c037d5686d364d7ae3a5469d1dfa}.
\end{document}
